\chapter{Certificación}
La Acreditación reconocía la competencia técnica de una organización para la realización de ciertas actividades, con las siguientes ventajas:
\begin{itemize}
    \item Declara que los organismos acreditados son competentes e imparciales;
    \item Les permite conseguir la aceptación de sus prestaciones a nivel internacional y el reconomiento de sus competencias.
    \item Evita a las empresas exportadoras los reiterados controles que deben pasar para tener acceso a los mercados internacionales
    \item Acceso a los mercados internacionales
    \item Establece y promueve la confianza a nivel nacional e internacional al comprobar la competencia de los operadores
\end{itemize}

\section{¿Qué es la certificación?}
La certificación es la acción llevada a cabo por una entidad reconocida como independiente de las partes interesadas, medianta la que se dispone de la confianza adecuada en que un producto, proceso o servicio debidamente identificado, es conforme con una norma u otro documento normativo especificado.

\section{¿Por qué se hace?}

El concepto de calidad ha evolucionado fuertemente a lo largo del siglo XX. se ha avanzado hacia la satisfacción de los requerimientos de los clientes, siguiendo luego por la adaptación para el costo, que implica incluir el aspecto económico.

La evolución de calidad tiene cuatro etapas:

\section{¿Qué se certifica?}
Actualmente podemos distinguir dos grandes grupos:

El HACCP es una herramienta para evaluar peligros y establecer sistemas de control centrados en la prevención. Puede aplicarse en toda la cadena alimentaria, desde el productor primario hasta el consumidor final.

\section{¿Cómo se hace?}

\subsection{Certificación de Sitemas de Calidad}

En esta categoría podemos encontrar la certificación de sistemas en base a las normas ISO 9000, ISO 14000, OSHAS 18000 u otras:
\begin{itemize}
    \item La ISO 9001:2000 y la ISO 14001:2004 dan los requisitos genéricos para los sistemas de gestión, no los requisitos para los productos específicos o los servicios.
    \item Las certificaciones de la ISO 9001:2000 y de la ISO 14001:2004 no son certificaciones del producto o garantías del producto.
    \item Se deben tomar las precauciones pertinentes de cualquier referencia a ella en la información relacionada con el producto, incluyendo los anuncios, o en cualquier otro medio, para evitar de dar la impresión que son certificaciones del producto o garantías del producto.
    \item En detalle, las marcas de certificación de la ISO 9001:2000 y de la ISO 14001:2004 NO deben ser exhibidas en rpoductos, en etiquetas del producto, en el producto que empaqueta, o de ninguna manera que se pueda interpretar como denotar conformidad del producto.
\end{itemize}

\subsection{Certificación de Productos}

\subsubsection{Modelos de Certificación de Producto ISO / CASCO}
Se describen los ocho modelos de Certificación existentes y la diferencia entre ellos, partiendo por el Ensayo de Tipo (Modelo Nº 1), pasando por la Inspección de Lote (Modelo Nº 7), para terminar en el Modelo Nº 8, que corresponde a la Inspección 100\%.

\paragraph{Modelo 1 - ISO CASCO 1}

\textbf{Ensayo de Tipo.} Servicio que certifica mediante ensayors técnicos de laboratorio y el análissi de los resultados obtenidos, la conformidad de una muestra de producto específica de acuerdo con lo establecido en normas técnicas, reglamentos técnicos o especificaciones técnicas pactadas entre el cliente y el proveedor.

La certificación de muestras no ampara otros lotes o muestras del mismo producto que no hayan sido evaluadas.